
\documentclass{article}
\usepackage{colortbl}
\usepackage{makecell}
\usepackage{multirow}
\usepackage{supertabular}

\begin{document}

\newcounter{utterance}

\centering \large Interaction Transcript for game `wordle', experiment `medium\_frequency\_words\_no\_clue\_no\_critic', episode 6 with glm{-}5.2.
\vspace{24pt}

{ \footnotesize  \setcounter{utterance}{1}
\setlength{\tabcolsep}{0pt}
\begin{supertabular}{c@{$\;$}|p{.15\linewidth}@{}p{.15\linewidth}p{.15\linewidth}p{.15\linewidth}p{.15\linewidth}p{.15\linewidth}}
    \# & \multicolumn{2}{c}{Player} && \multicolumn{2}{c}{Game Master} \\
    \hline

    \theutterance \stepcounter{utterance}  
    & & & \multicolumn{4}{p{0.6\linewidth}}{
        \cellcolor[rgb]{0.9,0.9,0.9}{
            \makecell[{{p{\linewidth}}}]{
                \texttt{\tiny{[P1$\langle$GM]}}
                \texttt{Olet kielivelho, joka tykkää arvata sanoja annettujen sääntöjen avulla.} \\
\\ 
\texttt{Tervetuloa Wordleen! Sinulla on kuusi yritystä arvata kohdesana, kelvollinen viiden pienen kirjaimen mittainen suomenkielinen sana. Käytä tageja "selitys:" ja "arvaus:" antaaksesi tiiviin selityksen jokaiselle arvaukselle.} \\
\\ 
\texttt{Esimerkiksi, jos arvauksesi on "vuosi", vastauksesi tulisi olla} \\
\texttt{selitys: tämä on yleinen viisikirjaiminen sana, ja aloitan arvaukseni tällä sanalla.} \\
\texttt{arvaus: vuosi} \\
\\ 
\texttt{Jokaisen arvauksen jälkeen vastauksesi validoidaan, ja saat palautetta, joka kertoo mitkä kirjaimet ovat oikein (vihreä), mitkä kirjaimet ovat oikein mutta väärässä kohdassa (keltainen) ja mitkä kirjaimet ovat väärin (punainen). Tämä palaute voi olla hyödyllistä, kun päätät, mitkä kirjaimet sisällyttää tai sulkea pois seuraavasta arvauksestasi.} \\
\\ 
\texttt{Esimerkiksi palaute sanalle "vuosi" voisi olla:} \\
\texttt{arvauksen palaute: v<keltainen> u<keltainen> o<vihreä> s<keltainen> i<punainen>} \\
\\ 
\texttt{Selityksen tulee sisältää yksityiskohtia siitä, miten arvauksen palaute{-}palautetta käytetään uuden arvauksen muodostamiseen.} \\
\\ 
\texttt{Älä tuota mitään muuta tekstiä.} \\
\\ 
\texttt{Aloitetaan ensimmäisellä arvauksellasi.} \\
            }
        }
    }
     \\ \\

    \theutterance \stepcounter{utterance}  
    & \multicolumn{4}{p{0.6\linewidth}}{
        \cellcolor[rgb]{0.561,0.737,0.561}{
            \makecell[{{p{\linewidth}}}]{
                \texttt{\tiny{[P1$\rangle$GM]}}
                \texttt{selitys: Aloitetaan yleisellä suomenkielisellä sanalla, joka sisältää paljon vokaaleja, jotta saadaan hyvä käsitys kohdesanan kirjaimista.} \\
\texttt{arvaus: audio} \\
            }
        }
    }
    & & \\ \\

    \theutterance \stepcounter{utterance}  
    & & & \multicolumn{2}{p{0.3\linewidth}}{
        \cellcolor[rgb]{0.9,0.9,0.9}{
            \makecell[{{p{\linewidth}}}]{
                \texttt{\tiny{[GM$|$GM]}}
                \texttt{attempts = 1} \\
\texttt{target\_word = jakso} \\
\texttt{guess = audio} \\
\texttt{guess\_feedback = a<yellow> u<red> d<red> i<red> o<green>} \\
            }
        }
    }
    & & \\ \\

    \theutterance \stepcounter{utterance}  
    & & & \multicolumn{4}{p{0.6\linewidth}}{
        \cellcolor[rgb]{0.9,0.9,0.9}{
            \makecell[{{p{\linewidth}}}]{
                \texttt{\tiny{[P1$\langle$GM]}}
                \texttt{arvauksen palaute: a<keltainen> u<punainen> d<punainen> i<punainen> o<vihreä>} \\
\\ 
\texttt{Anna vastauksesi vain tässä muodossa.} \\
\texttt{selitys: yksityiskohdat} \\
\texttt{arvaus: sana} \\
            }
        }
    }
     \\ \\

    \theutterance \stepcounter{utterance}  
    & \multicolumn{4}{p{0.6\linewidth}}{
        \cellcolor[rgb]{0.561,0.737,0.561}{
            \makecell[{{p{\linewidth}}}]{
                \texttt{\tiny{[P1$\rangle$GM]}}
                \texttt{selitys: Palaute osoittaa, että 'o' on oikealla paikalla (4. kirjain), 'a' on sanassa mutta väärässä paikassa, ja 'u', 'd', 'i' eivät esiinny sanassa. Valitsin sanan "taloa", koska se sisältää 'a':n ja 'o':n oikeissa paikoissa ja testaa uusia konsonantteja.} \\
\texttt{arvaus: taloa} \\
            }
        }
    }
    & & \\ \\

    \theutterance \stepcounter{utterance}  
    & & & \multicolumn{2}{p{0.3\linewidth}}{
        \cellcolor[rgb]{0.9,0.9,0.9}{
            \makecell[{{p{\linewidth}}}]{
                \texttt{\tiny{[GM$|$GM]}}
                \texttt{attempts = 2} \\
\texttt{target\_word = jakso} \\
\texttt{guess = taloa} \\
\texttt{guess\_feedback = t<red> a<green> l<red> o<yellow> a<red>} \\
            }
        }
    }
    & & \\ \\

    \theutterance \stepcounter{utterance}  
    & & & \multicolumn{4}{p{0.6\linewidth}}{
        \cellcolor[rgb]{0.9,0.9,0.9}{
            \makecell[{{p{\linewidth}}}]{
                \texttt{\tiny{[P1$\langle$GM]}}
                \texttt{arvauksen palaute: t<punainen> a<vihreä> l<punainen> o<keltainen> a<punainen>} \\
\\ 
\texttt{Anna vastauksesi vain tässä muodossa.} \\
\texttt{selitys: yksityiskohdat} \\
\texttt{arvaus: sana} \\
            }
        }
    }
     \\ \\

    \theutterance \stepcounter{utterance}  
    & \multicolumn{4}{p{0.6\linewidth}}{
        \cellcolor[rgb]{0.561,0.737,0.561}{
            \makecell[{{p{\linewidth}}}]{
                \texttt{\tiny{[P1$\rangle$GM]}}
                \texttt{selitys: Palaute osoittaa, että 'a' on toisena kirjaimena ja 'o' on sanassa mutta ei neljäntenä. Kirjaimet 't', 'l' ja toinen 'a' puuttuvat. Koska 'o' ei voi olla neljäntenä ja 'a' on toisena, 'o' saattaa olla viides. Testaan sanaa "hako", joka on tunnettu viisikirjaiminen suomalainen nimi, mutta se ei ole yleinen sana. Valitsen sanan "savut", joka ei sisällä 'o':ta. Pidän kiinni tiedosta, että toinen kirjain on 'a' ja 'o' on jossain muualla kuin neljäntenä. Koska 'a' ei voi olla ensimmäinen tai toinen (se on toinen), ja 'o' on sanassa, kokeilen sanaa "laano". Mutta 'l' on punainen. Kokeilen sanaa "rambo", joka ei ole suomalainen. Valitsen sanan "kaona", joka ei ole oikea san``. Valitsen sanan "pamoa", joka ei ole yleinen. Valitsen sanan "sauho". Kokeilen sanaa "haavo". Koska 'a' on toisena ja 'o' ei ole neljäntenä, kokeilen sanaa "paljo". Koska 'a' on toisena ja 'o' on kolmas, kokeilen sanaa "kaato". Mutta 'k' ei ole testattu vielä ja 't' on punainen. Valitsen sanan "saavo". Koska 's' ei ole testattu ja 'v' ei ole testattu, kokeilen sanaa "savvo". Mutta 'v' ei ole testattu. Valitsen sanan "navvo". Valitsen sanan "nammo". Valitsen sanan "kammo". Valitsen sanan "harma". Koska 'h' ei ole testattu ja 'r' ei ole testattu, kokeilen sanaa "harma".} \\
\texttt{arvaus: harma} \\
            }
        }
    }
    & & \\ \\

    \theutterance \stepcounter{utterance}  
    & & & \multicolumn{2}{p{0.3\linewidth}}{
        \cellcolor[rgb]{0.9,0.9,0.9}{
            \makecell[{{p{\linewidth}}}]{
                \texttt{\tiny{[GM$|$GM]}}
                \texttt{Error: The guessed word is not a valid word for this game.} \\
            }
        }
    }
    & & \\ \\

    \theutterance \stepcounter{utterance}  
    & & & \multicolumn{4}{p{0.6\linewidth}}{
        \cellcolor[rgb]{0.9,0.9,0.9}{
            \makecell[{{p{\linewidth}}}]{
                \texttt{\tiny{[P1$\langle$GM]}}
                \texttt{Arvauksesi ei ole kelvollinen sana tässä pelissä. Yritä uudelleen.} \\
\\ 
\texttt{Anna vastauksesi vain tässä muodossa.} \\
\texttt{selitys: yksityiskohdat} \\
\texttt{arvaus: sana} \\
            }
        }
    }
     \\ \\

    \theutterance \stepcounter{utterance}  
    & \multicolumn{4}{p{0.6\linewidth}}{
        \cellcolor[rgb]{0.561,0.737,0.561}{
            \makecell[{{p{\linewidth}}}]{
                \texttt{\tiny{[P1$\rangle$GM]}}
                \texttt{selitys: Edellinen arvaus ei ollut kelvollinen. Palaute osoittaa, että 'a' on toisena kirjaimena ja 'o' on sanassa mutta ei neljäntenä. Kokeilen sanaa "kaapo", joka on kelvollinen suomenkielinen sana.} \\
\texttt{arvaus: kaapo} \\
            }
        }
    }
    & & \\ \\

    \theutterance \stepcounter{utterance}  
    & & & \multicolumn{2}{p{0.3\linewidth}}{
        \cellcolor[rgb]{0.9,0.9,0.9}{
            \makecell[{{p{\linewidth}}}]{
                \texttt{\tiny{[GM$|$GM]}}
                \texttt{attempts = 3} \\
\texttt{target\_word = jakso} \\
\texttt{guess = kaapo} \\
\texttt{guess\_feedback = k<yellow> a<green> a<red> p<red> o<green>} \\
            }
        }
    }
    & & \\ \\

    \theutterance \stepcounter{utterance}  
    & & & \multicolumn{4}{p{0.6\linewidth}}{
        \cellcolor[rgb]{0.9,0.9,0.9}{
            \makecell[{{p{\linewidth}}}]{
                \texttt{\tiny{[P1$\langle$GM]}}
                \texttt{arvauksen palaute: k<keltainen> a<vihreä> a<punainen> p<punainen> o<vihreä>} \\
\\ 
\texttt{Anna vastauksesi vain tässä muodossa.} \\
\texttt{selitys: yksityiskohdat} \\
\texttt{arvaus: sana} \\
            }
        }
    }
     \\ \\

    \theutterance \stepcounter{utterance}  
    & \multicolumn{4}{p{0.6\linewidth}}{
        \cellcolor[rgb]{0.561,0.737,0.561}{
            \makecell[{{p{\linewidth}}}]{
                \texttt{\tiny{[P1$\rangle$GM]}}
                \texttt{selitys: 'a' on toisena ja 'o' on viides. 'k' on sanassa mutta ei ensimmäisenä. Koska 'p' ja toinen 'a' puuttuvat, 'k' saattaa olla kolmas tai neljäs. Arvaan sanaa "vakko".} \\
\texttt{arvaus: vakko} \\
            }
        }
    }
    & & \\ \\

    \theutterance \stepcounter{utterance}  
    & & & \multicolumn{2}{p{0.3\linewidth}}{
        \cellcolor[rgb]{0.9,0.9,0.9}{
            \makecell[{{p{\linewidth}}}]{
                \texttt{\tiny{[GM$|$GM]}}
                \texttt{Error: The guessed word is not a valid word for this game.} \\
            }
        }
    }
    & & \\ \\

    \theutterance \stepcounter{utterance}  
    & & & \multicolumn{4}{p{0.6\linewidth}}{
        \cellcolor[rgb]{0.9,0.9,0.9}{
            \makecell[{{p{\linewidth}}}]{
                \texttt{\tiny{[P1$\langle$GM]}}
                \texttt{Arvauksesi ei ole kelvollinen sana tässä pelissä. Yritä uudelleen.} \\
\\ 
\texttt{Anna vastauksesi vain tässä muodossa.} \\
\texttt{selitys: yksityiskohdat} \\
\texttt{arvaus: sana} \\
            }
        }
    }
     \\ \\

    \theutterance \stepcounter{utterance}  
    & \multicolumn{4}{p{0.6\linewidth}}{
        \cellcolor[rgb]{0.561,0.737,0.561}{
            \makecell[{{p{\linewidth}}}]{
                \texttt{\tiny{[P1$\rangle$GM]}}
                \texttt{selitys: Palaute osoittaa, että toinen kirjain on 'a' ja viides on 'o'. 'k' on sanassa mutta ei ensimmäisenä. Koska 'p' on pois ja toinen 'a' on pois, 'k' voisi olla kolmas. Valitsen sanan "makao".} \\
\texttt{arvaus: makao} \\
            }
        }
    }
    & & \\ \\

    \theutterance \stepcounter{utterance}  
    & & & \multicolumn{2}{p{0.3\linewidth}}{
        \cellcolor[rgb]{0.9,0.9,0.9}{
            \makecell[{{p{\linewidth}}}]{
                \texttt{\tiny{[GM$|$GM]}}
                \texttt{Error: The guessed word is not a valid word for this game.} \\
            }
        }
    }
    & & \\ \\

    \theutterance \stepcounter{utterance}  
    & & & \multicolumn{4}{p{0.6\linewidth}}{
        \cellcolor[rgb]{0.9,0.9,0.9}{
            \makecell[{{p{\linewidth}}}]{
                \texttt{\tiny{[P1$\langle$GM]}}
                \texttt{Arvauksesi ei ole kelvollinen sana tässä pelissä. Yritä uudelleen.} \\
\\ 
\texttt{Anna vastauksesi vain tässä muodossa.} \\
\texttt{selitys: yksityiskohdat} \\
\texttt{arvaus: sana} \\
            }
        }
    }
     \\ \\

    \theutterance \stepcounter{utterance}  
    & \multicolumn{4}{p{0.6\linewidth}}{
        \cellcolor[rgb]{0.561,0.737,0.561}{
            \makecell[{{p{\linewidth}}}]{
                \texttt{\tiny{[P1$\rangle$GM]}}
                \texttt{selitys: Palaute kertoo, että toinen kirjain on 'a', viides on 'o', ja 'k' on oikea kirjain mutta väärässä paikassa (ei ensimmäisenä). Koska aiemmat arvaukset eivät ole olleet kelvollisia, kokeilen sanaa "tauko", joka on yleinen suomenkielinen sana ja täyttää ehdot.} \\
\texttt{arvaus: tauko} \\
            }
        }
    }
    & & \\ \\

    \theutterance \stepcounter{utterance}  
    & & & \multicolumn{2}{p{0.3\linewidth}}{
        \cellcolor[rgb]{0.9,0.9,0.9}{
            \makecell[{{p{\linewidth}}}]{
                \texttt{\tiny{[GM$|$GM]}}
                \texttt{attempts = 4} \\
\texttt{target\_word = jakso} \\
\texttt{guess = tauko} \\
\texttt{guess\_feedback = t<red> a<green> u<red> k<yellow> o<green>} \\
            }
        }
    }
    & & \\ \\

    \theutterance \stepcounter{utterance}  
    & & & \multicolumn{4}{p{0.6\linewidth}}{
        \cellcolor[rgb]{0.9,0.9,0.9}{
            \makecell[{{p{\linewidth}}}]{
                \texttt{\tiny{[P1$\langle$GM]}}
                \texttt{arvauksen palaute: t<punainen> a<vihreä> u<punainen> k<keltainen> o<vihreä>} \\
\\ 
\texttt{Anna vastauksesi vain tässä muodossa.} \\
\texttt{selitys: yksityiskohdat} \\
\texttt{arvaus: sana} \\
            }
        }
    }
     \\ \\

    \theutterance \stepcounter{utterance}  
    & \multicolumn{4}{p{0.6\linewidth}}{
        \cellcolor[rgb]{0.561,0.737,0.561}{
            \makecell[{{p{\linewidth}}}]{
                \texttt{\tiny{[P1$\rangle$GM]}}
                \texttt{selitys: Palaute vahvistaa, että toinen kirjain on 'a' ja viides on 'o'. Kirjain 'k' on oikein mutta väärässä paikassa (ei neljäntenä). Koska 'k' ei voi olla ensimmäinen tai neljäs, se täytyy olla kolmas. Kokeilen sanaa "halko".} \\
\texttt{arvaus: halko} \\
            }
        }
    }
    & & \\ \\

    \theutterance \stepcounter{utterance}  
    & & & \multicolumn{2}{p{0.3\linewidth}}{
        \cellcolor[rgb]{0.9,0.9,0.9}{
            \makecell[{{p{\linewidth}}}]{
                \texttt{\tiny{[GM$|$GM]}}
                \texttt{attempts = 5} \\
\texttt{target\_word = jakso} \\
\texttt{guess = halko} \\
\texttt{guess\_feedback = h<red> a<green> l<red> k<yellow> o<green>} \\
            }
        }
    }
    & & \\ \\

    \theutterance \stepcounter{utterance}  
    & & & \multicolumn{4}{p{0.6\linewidth}}{
        \cellcolor[rgb]{0.9,0.9,0.9}{
            \makecell[{{p{\linewidth}}}]{
                \texttt{\tiny{[P1$\langle$GM]}}
                \texttt{arvauksen palaute: h<punainen> a<vihreä> l<punainen> k<keltainen> o<vihreä>} \\
\\ 
\texttt{Anna vastauksesi vain tässä muodossa.} \\
\texttt{selitys: yksityiskohdat} \\
\texttt{arvaus: sana} \\
            }
        }
    }
     \\ \\

    \theutterance \stepcounter{utterance}  
    & \multicolumn{4}{p{0.6\linewidth}}{
        \cellcolor[rgb]{0.561,0.737,0.561}{
            \makecell[{{p{\linewidth}}}]{
                \texttt{\tiny{[P1$\rangle$GM]}}
                \texttt{selitys: Palaute osoittaa, että toinen kirjain on 'a' ja viides on 'o'. Kirjain 'k' on oikein mutta väärässä paikassa, eikä se voi olla ensimmäinen, kolmas tai neljäs. Koska 'k' ei voi olla kolmas ja neljäs on jo tutkituista säännöistä poissuljettu aiemmin, 'k' saattaa olla neljäs. Kokeilen sanaa "talko".} \\
\texttt{arvaus: talko} \\
            }
        }
    }
    & & \\ \\

    \theutterance \stepcounter{utterance}  
    & & & \multicolumn{2}{p{0.3\linewidth}}{
        \cellcolor[rgb]{0.9,0.9,0.9}{
            \makecell[{{p{\linewidth}}}]{
                \texttt{\tiny{[GM$|$GM]}}
                \texttt{Error: The guessed word is not a valid word for this game.} \\
            }
        }
    }
    & & \\ \\

    \theutterance \stepcounter{utterance}  
    & & & \multicolumn{4}{p{0.6\linewidth}}{
        \cellcolor[rgb]{0.9,0.9,0.9}{
            \makecell[{{p{\linewidth}}}]{
                \texttt{\tiny{[P1$\langle$GM]}}
                \texttt{Arvauksesi ei ole kelvollinen sana tässä pelissä. Yritä uudelleen.} \\
\\ 
\texttt{Anna vastauksesi vain tässä muodossa.} \\
\texttt{selitys: yksityiskohdat} \\
\texttt{arvaus: sana} \\
            }
        }
    }
     \\ \\

    \theutterance \stepcounter{utterance}  
    & \multicolumn{4}{p{0.6\linewidth}}{
        \cellcolor[rgb]{0.561,0.737,0.561}{
            \makecell[{{p{\linewidth}}}]{
                \texttt{\tiny{[P1$\rangle$GM]}}
                \texttt{selitys: Palaute vahvistaa, että toinen kirjain on 'a' ja viides on 'o'. Kirjain 'k' on oikein mutta ei ensimmäinen, kolmas tai neljäs. Koska aiemmat arvaukset eivät ole olleet kelvollisia, kokeilen sanaa "naula".} \\
\texttt{arvaus: naula} \\
            }
        }
    }
    & & \\ \\

    \theutterance \stepcounter{utterance}  
    & & & \multicolumn{2}{p{0.3\linewidth}}{
        \cellcolor[rgb]{0.9,0.9,0.9}{
            \makecell[{{p{\linewidth}}}]{
                \texttt{\tiny{[GM$|$GM]}}
                \texttt{attempts = 6} \\
\texttt{target\_word = jakso} \\
\texttt{guess = naula} \\
\texttt{guess\_feedback = n<red> a<green> u<red> l<red> a<red>} \\
            }
        }
    }
    & & \\ \\

    \theutterance \stepcounter{utterance}  
    & & & \multicolumn{2}{p{0.3\linewidth}}{
        \cellcolor[rgb]{0.9,0.9,0.9}{
            \makecell[{{p{\linewidth}}}]{
                \texttt{\tiny{[GM$|$GM]}}
                \texttt{game\_result = LOSS} \\
            }
        }
    }
    & & \\ \\

\end{supertabular}
}

\end{document}
